% topic_04.tex — Content only. Compiled via main.tex.
% ── No \documentclass, \begin{document} or \end{document} here ──

\chapteropener{4}{Forces, Density and Pressure}{%
  \item Turning Effects of Forces
  \item Equilibrium of Forces
  \item Density and Pressure
  \item Hydrostatic Pressure
  \item Upthrust and Archimedes' Principle
}

% ===== SECTION 1: TURNING EFFECTS =====
\section{Turning Effects of Forces}

\begin{definitionbox}{Centre of Gravity}
The \textbf{centre of gravity} of an object is the single point at which the entire weight of the object may be considered to act.
\begin{itemize}[itemsep=2pt]
  \item For a uniform, regular object the centre of gravity lies at the geometric centre.
  \item For irregular objects it can be found experimentally by suspension.
  \item All gravitational calculations treat the object as a \textbf{point mass} at this location.
\end{itemize}
\end{definitionbox}

\begin{definitionbox}{Moment of a Force}
The \textbf{moment} of a force about a point is the product of the force and the perpendicular distance from the point to the line of action of the force.
$$\text{moment} = F \times d_\perp \qquad \text{units: N\,m}$$
where $d_\perp$ is the \textbf{perpendicular} distance from the pivot to the line of action.
\end{definitionbox}

\begin{center}
\begin{tikzpicture}[scale=0.9]
\usetikzlibrary{angles,quotes}
\begin{scope}[xshift=0cm]
  % Beam
  \draw[darkblue, very thick] (-3,0) -- (3,0);
  % Pivot triangle
  \fill[darkblue!60] (-0.25,-0.5) -- (0.25,-0.5) -- (0,-0.1) -- cycle;
  \draw[darkgray, thick] (-0.6,-0.5) -- (0.6,-0.5);
  % Force arrow
  \draw[-latex, midblue, very thick] (2.2,1.5) -- (2.2,0.05);
  \node[midblue, font=\small\bfseries, anchor=west] at (2.35,0.9) {$F$};
  % Perpendicular distance
  \draw[latex-latex, darkgray, thick] (0,-0.85) -- (2.2,-0.85);
  \node[darkgray, font=\small] at (1.1,-1.15) {$d_\perp$};
  % Dashed line from pivot to force
  \draw[dashed, darkgray] (0,0) -- (0,-0.85);
  \draw[dashed, darkgray] (2.2,0) -- (2.2,-0.85);
  \node[darkblue, font=\small\bfseries] at (0,-1.6) {Moment $= F \times d_\perp$};
  \end{scope}

\begin{scope}[xshift=4cm,scale = 0.8]

   \coordinate (A) at (8,0);
    \coordinate (O) at (6,0);
    \coordinate (B) at (8,2);
  % Beam
  \draw[darkblue, very thick] (0,0) -- (6,0);
  \draw[dashed, darkgray] (6,0) -- (9,0);
 
  % Force arrow
  \draw[-latex, midblue, very thick] (6,0) -- (8,2);
  \node[midblue, font=\small\bfseries, anchor=west] at (5.7,1.1) {$F$};
  
   % distance
  \draw[latex-latex, darkgray, thick] (0,0.2) -- (6,0.2);
  \node[darkgray, font=\small] at (3,0.5) {$d$};
  
  % Perpendicular distance
  \draw[latex-latex, darkgray, thick] (0,0) -- (3,-3);
  \node[darkgray, font=\small,fill=white] at (1.1,-1.8) {$d_\perp=dsin\theta$};
  % Dashed line from pivot to force
  \draw[dashed, darkgray] (6,0) -- (3,-3);
 % \draw[dashed, darkgray] (0,0) -- (3,-3);
  %\node[darkblue, font=\small\bfseries] at (0,-1.6) {Moment $= F \times d_\perp$};
   \pic [draw, ->, "$\theta$", angle radius=1cm, angle eccentricity=1.5] 
        {angle = A--O--B};
  \end{scope}
\end{tikzpicture}
\end{center}

\begin{warningbox}{Common Mistake}
Always use the \textbf{perpendicular} distance from the pivot to the \emph{line of action} of the force — not the distance along the beam or object. If the force is at an angle $\theta$, the moment is $F\sin\theta \times d$ or equivalently $F \times d\sin\theta$.
\end{warningbox}

\pagebreak

\subsection{Couples and Torque}

\begin{definitionbox}{Couple}
A \textbf{couple} is a pair of forces that:
\begin{itemize}[itemsep=2pt]
  \item are \textbf{equal in magnitude},
  \item act in \textbf{opposite directions},
  \item have \textbf{parallel but different lines of action}.
\end{itemize}
A couple produces \textbf{rotation only} — it has zero resultant force, so it produces no translational acceleration.
\end{definitionbox}

\begin{formulabox}{Torque of a Couple}
\begin{center}
\Large $\tau = F \times d$
\end{center}
\vspace{0.3em}
\begin{tabular}{@{}ll@{}}
$\tau$ & = torque of the couple (N\,m) \\
$F$   & = magnitude of one of the forces (N) \\
$d$   & = perpendicular distance between the lines of action of the forces (m) \\
\end{tabular}

\vspace{0.5em}
Note: unlike the moment of a single force, the torque of a couple is the \textbf{same about any point}.
\end{formulabox}

\vspace{4em}

\begin{center}
\begin{tikzpicture}[scale=1.6]
  % Beam
  \draw[darkblue, very thick] (-2.5,0) -- (2.5,0);
  % Top force (left)
  \draw[->, midblue, very thick] (-2,0.1) -- (-2,1.5);
  \node[midblue, font=\small\bfseries, anchor=east] at (-2.15,0.9) {$F$};
  % Bottom force (right)
  \draw[->, midblue, very thick] (2,-0.1) -- (2,-1.5);
  \node[midblue, font=\small\bfseries, anchor=west] at (2.15,-0.9) {$F$};
  % Distance arrow
  \draw[<->, darkgray, thick] (-2,-1.9) -- (2,-1.9);
  \node[darkgray, font=\small] at (0,-2.2) {$d$};
  \draw[dashed, darkgray] (-2,-0.1) -- (-2,-1.9);
  \draw[dashed, darkgray] (2,-0.1) -- (2,-1.9);
  % Label
  \node[darkblue, font=\small\bfseries] at (0,1.7) {Torque $\tau = F \times d$};
\end{tikzpicture}
\end{center}

\pagebreak

% ===== SECTION 2: EQUILIBRIUM =====
\section{Equilibrium of Forces}

\begin{definitionbox}{Conditions for Equilibrium}
A body is in \textbf{equilibrium} when:
\begin{enumerate}[itemsep=2pt]
  \item The \textbf{resultant force} is zero (no translational acceleration).
  \item The \textbf{resultant torque} about any point is zero (no rotational acceleration).
\end{enumerate}
Both conditions must be satisfied simultaneously.
\end{definitionbox}

\begin{formulabox}{Principle of Moments}
For a body in equilibrium:
\begin{center}
\large $\sum \text{clockwise moments} = \sum \text{anticlockwise moments}$
\end{center}
(about any chosen pivot point)
\end{formulabox}

\begin{keybox}{Using the Principle of Moments}
\begin{itemize}[itemsep=2pt]
  \item Choose a \textbf{convenient pivot} — often at the point of an unknown force, so that unknown disappears from the moment equation.
  \item List all forces and their perpendicular distances from the chosen pivot.
  \item Apply the principle; then use $\sum F = 0$ to find any remaining unknowns.
  \item Check both conditions of equilibrium are satisfied.
\end{itemize}
\end{keybox}



\subsection{Coplanar Forces in Equilibrium — Vector Triangle}

\begin{keybox}{Three Coplanar Forces in Equilibrium}
If exactly three coplanar forces act on a body in equilibrium, they must be \textbf{concurrent} (pass through a single point) and can be represented by the three sides of a \textbf{closed triangle} drawn head-to-tail.
\begin{itemize}[itemsep=2pt]
  \item Draw the force vectors end-to-end.
  \item The triangle must close — meaning the resultant is zero.
  \item Trigonometry (sine rule, cosine rule, or right-triangle ratios) can then be used to find unknown magnitudes or directions.
\end{itemize}
\end{keybox}

\begin{center}
\begin{tikzpicture}[scale=.9]
 \begin{scope}[xshift=0cm]
  % Three force coplanar forces
  \coordinate (A) at (0,0);
  \coordinate (B) at (3,1.5);
  \coordinate (C) at (-3,1.5);
  \coordinate (D) at (0,-3);
  \draw[-latex, midblue, very thick] (A) -- (B) node[midway, below, font=\small] {$F_2$};
  \draw[-latex, darkblue, very thick] (A) -- (C) node[midway, below, font=\small] {$F_1$};
 \draw[-latex, pinkframe, very thick] (A) -- (D) node[midway, left, font=\small] {$F_3$};
  %\node[darkgray, font=\small\itshape] at (1.4,-0.55) {Closed triangle $\Rightarrow$ equilibrium};
  \end{scope}



\begin{scope}[xshift=8cm, yshift=-2cm]
 % Three force vectors forming a closed triangle
  \coordinate (A) at (0,0);
  \coordinate (B) at (-3,1.5);
  \coordinate (C) at (0,3);
 
  \draw[-latex, darkblue, very thick] (A) -- (B) node[midway, below, font=\small] {$F_1$};
  \draw[-latex, midblue, very thick] (B) -- (C) node[midway, above, font=\small] {$F_2$};
 \draw[-latex, pinkframe, very thick] (C) -- (A) node[midway, left, font=\small] {$F_3$};
  \node[darkgray, font=\small\itshape] at (-1,-0.55) {Closed triangle $\Rightarrow$ equilibrium};
  
  \end{scope}
\end{tikzpicture}
\end{center}

% ===== SECTION 3: DENSITY AND PRESSURE =====
\section{Density and Pressure}

\begin{definitionbox}{Density}
\textbf{Density} is defined as mass per unit volume:
$$\rho = \frac{m}{V} \qquad \text{units: kg\,m}^{-3}$$
\begin{tabular}{@{}ll@{}}
$\rho$ & = density (kg\,m$^{-3}$) \\
$m$   & = mass (kg) \\
$V$   & = volume (m$^3$) \\
\end{tabular}
\end{definitionbox}

\begin{definitionbox}{Pressure}
\textbf{Pressure} is defined as the normal force per unit area:
$$P = \frac{F}{A} \qquad \text{units: Pa} \equiv \text{N\,m}^{-2}$$
\begin{tabular}{@{}ll@{}}
$P$ & = pressure (Pa) \\
$F$ & = normal force acting on the surface (N) \\
$A$ & = area over which the force acts (m$^2$) \\
\end{tabular}
\end{definitionbox}

\begin{warningbox}{Pressure is a Scalar}
Although pressure arises from a force (a vector), pressure itself is a \textbf{scalar} quantity — it acts equally in all directions at a point within a fluid. Do not give pressure a direction in your answers.
\end{warningbox}


\pagebreak
% ===== SECTION 4: HYDROSTATIC PRESSURE =====
\section{Hydrostatic Pressure}


\begin{center}
% WATER COLUMN WEIGHT
\colorlet{mydarkblue}{blue!50!black}
\colorlet{myred}{red!65!black}
\colorlet{vcol}{green!45!black}
\colorlet{watercol}{blue!80!cyan!10!white}
\colorlet{darkwatercol}{blue!80!cyan!20!white}
\colorlet{metalcol}{blue!40!black!10!white}
\tikzstyle{water}=[draw=mydarkblue,top color=watercol!90,bottom color=watercol!90!black,shading angle=5]
\tikzstyle{vertical water}=[water,
  top color=watercol!90!black!90,bottom color=watercol!90!black!90,middle color=watercol!80,shading angle=90]
\tikzstyle{dark water}=[draw=mydarkblue,top color=darkwatercol,bottom color=darkwatercol!80!black,shading angle=5]
\begin{tikzpicture}
  \def\Rx{1.1}      % tank horizontal radius
  \def\Ry{0.5}      % tank vertical radius
  \def\rx{0.22*\Rx} % column horizontal radius
  \def\ry{0.22*\Ry} % column vertical radius
  \def\H{2.6}       % height tank
  \def\h{0.94*\H}   % height water
  \def\y{0.40*\H}   % vertical position piece
  \def\dy{0.10*\H}  % height piece
  
  % WATER
  \draw[vertical water]
    (-\Rx,\h) -- (-\Rx,0) arc (180:360:{\Rx} and {\Ry}) -- (\Rx,\h);
  \draw[water]
    (0,\h) ellipse ({\Rx} and {\Ry});
  
  % COLUMN + PIECE
  \draw[dark water]
    (-\rx,\y) -- (-\rx,\y-\dy) arc (180:360:{\rx} and {\ry}) -- (\rx,\y);
  \draw[dark water]
    (0,\y) ellipse ({\rx} and {\ry});
  \draw[dark water,opacity=0.50,draw=none] % column
    (-\rx,\h) -- (-\rx,\y) arc (180:360:{\rx} and {\ry}) -- (\rx,\h) arc(360:180:{\rx} and {\ry}) -- cycle;
  \draw[blue!20!black,dashed,very thin,opacity=0.50] (-\rx,\y) -- (-\rx,\h) (\rx,\y) -- (\rx,\h);
  \draw[dark water,opacity=0.50] % top column
    (0,\h) ellipse ({\rx} and {\ry});
  \draw[blue!30!black]
    (0,\h) ellipse ({\Rx} and {\Ry});
  
  % CONTAINER
  \draw[thick]
    (-\Rx,\H) -- (-\Rx,0) arc (180:360:{\Rx} and {\Ry}) -- (\Rx,\H);
  \draw[thick]
    (0,\H) ellipse ({\Rx} and {\Ry});
  
  % LABELS
  \node[scale=0.98] at ({-0.45*(\Rx+\rx)},1.04*\h) {$P_\mathrm{atm}$}; %{\contour{watercol}{$P_\mathrm{atm}$}};
  \node at ({-0.40*(\Rx+\rx)},{\y-0.5*\dy}) {$P$};
  \node[above=-2] at (0,\y+\ry) {$A$};
  \draw[<->] (1.4*\Rx,\y) --++ (0,\h-\y) node[midway,fill=white,inner sep=1] {$\Delta h$};
  
\end{tikzpicture}
\end{center}

\begin{formulabox}{Hydrostatic Pressure}

The additional pressure on a submerged object is due to the weight of the fluid column above it. In the diagram this is a cylinder of Volume $A\Delta h$

\begin{itemize}[itemsep=4pt]
  \item Weight of column of fluid: $W = mg = \rho V g = \rho \times (A\,\Delta h) \times g$
  \item Extra pressure on an Area $A$: $\Delta P = \dfrac{W}{A} = \dfrac{\rho A\,\Delta h\,g}{A}$
\end{itemize}

\begin{center}
\Large $\Delta P = \rho g \Delta h$
\end{center}
\vspace{0.3em}
\begin{tabular}{@{}ll@{}}
$\Delta P$  & = increase in pressure with depth (Pa) \\
$\rho$      & = density of the fluid (kg\,m$^{-3}$) \\
$g$         & = gravitational field strength (N\,kg$^{-1}$) \\
$\Delta h$  & = increase in depth below the surface (m) \\
\end{tabular}
\end{formulabox}

\begin{keybox}{Key Points about Hydrostatic Pressure}
\begin{itemize}[itemsep=2pt]
  \item Pressure depends only on \textbf{depth} $\Delta h$, not on the shape or cross-sectional area of the container.
  \item Pressure acts \textbf{equally in all directions} at a given depth.
  \item In connected vessels, fluid reaches the same height regardless of the vessel shape (\emph{communicating vessels}).
  \item Total absolute pressure at depth $h$: $P = P_\mathrm{atm} + \rho g h$ 
\end{itemize}
\end{keybox}


\pagebreak

% ===== SECTION 5: UPTHRUST AND ARCHIMEDES =====
\section{Upthrust and Archimedes' Principle}

\begin{definitionbox}{Origin of Upthrust}
\textbf{Upthrust} (buoyancy force) acts on any object immersed in a fluid. It arises because the \textbf{hydrostatic pressure on the bottom face} of the object is greater than on the top face (since the bottom is at greater depth). The net upward pressure force is the upthrust.
\end{definitionbox}

\begin{formulabox}{Archimedes' Principle}
\begin{center}
\Large $F = \rho g V$
\end{center}
\vspace{0.3em}
\begin{tabular}{@{}ll@{}}
$F$   & = upthrust (N) \\
$\rho$& = density of the fluid (kg\,m$^{-3}$) \\
$g$   & = gravitational field strength (N\,kg$^{-1}$) \\
$V$   & = volume of fluid displaced by the object (m$^3$) \\
\end{tabular}

\vspace{0.5em}
\textit{Archimedes' Principle:} the upthrust equals the \textbf{weight of fluid displaced}.
\end{formulabox}

\begin{center}
\begin{tikzpicture}[scale=0.85]
  % Fluid container
  \draw[darkblue, thick] (-2,0) -- (-2,3.8) -- (2,3.8) -- (2,0) -- cycle;
  \fill[midblue!20] (-1.95,0.05) rectangle (1.95,3.75);
  % Submerged block
  \fill[darkblue!60] (-0.7,0.9) rectangle (0.7,2.1);
  \draw[darkblue, thick] (-0.7,0.9) rectangle (0.7,2.1);
  % Weight arrow
  \draw[-latex, pinkframe, very thick] (0,1.5) -- (0,0.3);
  \node[pinkframe, font=\small\bfseries, anchor=west] at (0,0.6) {$W = mg$};
  % Upthrust arrow
  \draw[-latex, midblue, very thick] (0,1.5) -- (0,2.9);
  \node[midblue, font=\small\bfseries, anchor=west] at (0,2.5) {$F = \rho g V$};
  % Pressure labels
  \node[midblue, font=\scriptsize] at (-1.3,0.9) {high $p$};
  \node[midblue, font=\scriptsize] at (-1.3,2.1) {low $p$};
\end{tikzpicture}
\end{center}

\begin{keybox}{Floating and Sinking}
\begin{itemize}[itemsep=2pt]
  \item If $F > W$: net upward force — object \textbf{rises}.
  \item If $F = W$: object is in equilibrium — \textbf{floats} (fully or partially submerged).
  \item If $F < W$: net downward force — object \textbf{sinks} (rests on the bottom).
  \item A floating object displaces fluid whose weight \textbf{equals} the object's weight.
\end{itemize}
\end{keybox}


\pagebreak


% ===== SECTION 6: WORKED EXAMPLES =====
\section{Worked Examples}

\subsection*{Example 1 — Principle of Moments}

\textbf{Question:} A uniform beam of weight $120\ \text{N}$ and length $2.0\ \text{m}$ is pivoted at one end. A $200\ \text{N}$ load hangs $1.5\ \text{m}$ from the pivot. Calculate the vertical force $F$ needed at the free end to maintain equilibrium.

\begin{examplebox}{Solution}
\textbf{Solution:}\\
Take moments about the pivot (left end):

Clockwise moments: weight of beam acts at centre $(1.0\ \text{m})$ and load at $1.5\ \text{m}$:
$$\sum M_\text{CW} = (120 \times 1.0) + (200 \times 1.5) = 120 + 300 = 420\ \text{N\,m}$$

Anticlockwise moment from $F$ at $2.0\ \text{m}$:
$$\sum M_\text{ACW} = F \times 2.0$$

Setting equal: $F \times 2.0 = 420$, so $\boxed{F = 210\ \text{N}}$ (upwards).
\end{examplebox}

\subsection*{Example 2 — Hydrostatic Pressure}

\textbf{Question:} A submarine is at a depth of $250\ \text{m}$ in seawater of density $1025\ \text{kg\,m}^{-3}$. Calculate the additional pressure at this depth compared with the surface. ($g = 9.81\ \text{N\,kg}^{-1}$)

\begin{examplebox}{Solution}
\textbf{Solution:}\\
$$\Delta p = \rho g \Delta h = 1025 \times 9.81 \times 250$$
$$\Delta p = 1025 \times 9.81 \times 250 = \mathbf{2.51 \times 10^6\ \text{Pa}}\ (= 2.51\ \text{MPa})$$

This is roughly 25 times atmospheric pressure — illustrating the enormous pressures at depth.
\end{examplebox}

\subsection*{Example 3 — Upthrust (Archimedes' Principle)}

\textbf{Question:} A solid aluminium sphere of radius $0.080\ \text{m}$ is fully submerged in water (density $1000\ \text{kg\,m}^{-3}$). Calculate the upthrust acting on the sphere.

\begin{examplebox}{Solution}
\textbf{Solution:}\\
Volume of sphere: $V = \dfrac{4}{3}\pi r^3 = \dfrac{4}{3}\pi (0.080)^3 = 2.14 \times 10^{-3}\ \text{m}^3$

$$F = \rho g V = 1000 \times 9.81 \times 2.14\times10^{-3} = \mathbf{21.0\ \text{N}}$$

Density of aluminium $\approx 2700\ \text{kg\,m}^{-3}$, so its weight $= 2700 \times 9.81 \times 2.14\times10^{-3} \approx 56.6\ \text{N}$. Since $W > F$, the sphere sinks.
\end{examplebox}

% ===== SECTION 7: PRACTICE QUESTIONS =====
\section{Practice Exam Questions}

\subsection*{Section A — Short Answer Questions}

\begin{tcolorbox}[colback=white, colframe=midblue, breakable]

\begin{minipage}{\linewidth}
\textbf{Q1.} Define the moment of a force and state its SI unit.
\par\hfill\textit{[2 marks]}
\vspace{2.5cm}
\end{minipage}

\begin{minipage}{\linewidth}
\textbf{Q2.} State the two conditions that must be satisfied for a rigid body to be in equilibrium.
\par\hfill\textit{[2 marks]}
\vspace{2.5cm}
\end{minipage}

\begin{minipage}{\linewidth}
\textbf{Q3.} Define density and pressure, giving the SI unit of each.
\par\hfill\textit{[4 marks]}
\vspace{3cm}
\end{minipage}

\begin{minipage}{\linewidth}
\textbf{Q4.} Explain, using the concept of pressure, why upthrust acts on a submerged object.
\par\hfill\textit{[2 marks]}
\vspace{3cm}
\end{minipage}

\begin{minipage}{\linewidth}
\textbf{Q5.} A couple consists of two forces, each of magnitude $35\ \text{N}$, separated by a perpendicular distance of $0.24\ \text{m}$. Calculate the torque of the couple.
\par\hfill\textit{[2 marks]}
\vspace{2.5cm}
\end{minipage}

\end{tcolorbox}

\subsection*{Section B — Longer Structured Questions}

\begin{tcolorbox}[colback=white, colframe=midblue, breakable]

\begin{minipage}{\linewidth}
\textbf{Q6.} A uniform plank of mass $8.0\ \text{kg}$ and length $3.0\ \text{m}$ is supported at each end by vertical forces $F_1$ (at the left) and $F_2$ (at the right). A person of mass $60\ \text{kg}$ stands $1.0\ \text{m}$ from the left end.
\end{minipage}

\begin{enumerate}[label=(\alph*)]
  \item \begin{minipage}[t]{\linewidth}
  Calculate the force $F_2$ at the right support.
  \par\hfill\textit{[3 marks]}
  \vspace{4cm}
  \end{minipage}

  \item \begin{minipage}[t]{\linewidth}
  Calculate the force $F_1$ at the left support.
  \par\hfill\textit{[2 marks]}
  \vspace{3cm}
  \end{minipage}

  \item \begin{minipage}[t]{\linewidth}
  The person moves towards the right end. Describe and explain what happens to $F_1$ and $F_2$.
  \par\hfill\textit{[2 marks]}
  \vspace{3.5cm}
  \end{minipage}
\end{enumerate}

\begin{minipage}{\linewidth}
\textbf{Q7.} A diving bell (a sealed metal container open at the bottom) is lowered into the sea. The interior of the bell initially contains air at atmospheric pressure $p_0 = 1.01 \times 10^5\ \text{Pa}$.
\end{minipage}

\begin{enumerate}[label=(\alph*)]
  \item \begin{minipage}[t]{\linewidth}
  Derive the equation $\Delta p = \rho g \Delta h$ for hydrostatic pressure, starting from the definitions of pressure and density.
  \par\hfill\textit{[3 marks]}
  \vspace{5cm}
  \end{minipage}

  \item \begin{minipage}[t]{\linewidth}
  The bell is lowered to a depth of $180\ \text{m}$ in seawater of density $1025\ \text{kg\,m}^{-3}$. Calculate the total pressure at this depth.
  \par\hfill\textit{[2 marks]}
  \vspace{3.5cm}
  \end{minipage}

  \item \begin{minipage}[t]{\linewidth}
  A steel sphere of volume $4.5 \times 10^{-3}\ \text{m}^3$ and mass $35\ \text{kg}$ is attached beneath the bell. Calculate the upthrust on the sphere and determine whether it floats or sinks when released.
  \par\hfill\textit{[3 marks]}
  \vspace{4.5cm}
  \end{minipage}
\end{enumerate}

\end{tcolorbox}

\pagebreak


% ===== SECTION 8: MARK SCHEME =====
\section{Mark Scheme and Answers}

\begin{markscheme}

\textbf{Q1.} The moment of a force is the product of the force and the perpendicular distance from the pivot to the line of action of the force [1]; unit: N\,m [1].

\vspace{0.5em}\textbf{Q2.} The resultant force on the body is zero [1]; and the resultant torque (moment) about any point is zero [1].

\vspace{0.5em}\textbf{Q3.} Density: mass per unit volume [1], kg\,m$^{-3}$ [1]. Pressure: (normal) force per unit area [1], Pa (or N\,m$^{-2}$) [1].

\vspace{0.5em}\textbf{Q4.} The pressure increases with depth [1]; so the upward pressure force on the bottom face of the object is greater than the downward pressure force on the top face, giving a net upward force [1].

\vspace{0.5em}\textbf{Q5.} $\tau = Fd = 35 \times 0.24 = \mathbf{8.4\ \text{N\,m}}$ [2].

\vspace{0.5em}\textbf{Q6(a).} Take moments about left end. Clockwise: weight of plank $= 8.0 \times 9.81 = 78.5\ \text{N}$ at $1.5\ \text{m}$; person's weight $= 60 \times 9.81 = 589\ \text{N}$ at $1.0\ \text{m}$. Anticlockwise: $F_2$ at $3.0\ \text{m}$. Principle of moments: $F_2 \times 3.0 = (589 \times 1.0) + (78.5 \times 1.5)$ [1]; $F_2 \times 3.0 = 589 + 117.8 = 706.8$ [1]; $F_2 = \mathbf{236\ \text{N}}$ [1].

\vspace{0.5em}\textbf{Q6(b).} $\sum F = 0$: $F_1 + F_2 = 589 + 78.5 = 667.5\ \text{N}$ [1]; $F_1 = 667.5 - 236 = \mathbf{432\ \text{N}}$ [1].

\vspace{0.5em}\textbf{Q6(c).} As the person moves right, their moment about the left end increases, so $F_2$ increases [1]; and since the total must remain constant, $F_1$ decreases [1].

\vspace{0.5em}\textbf{Q7(a).} Consider a horizontal slab of fluid, area $A$, depth $\Delta h$, density $\rho$ [1]. Weight of slab $= \rho A\,\Delta h\,g$ [1]; extra pressure $= W/A = \rho g\,\Delta h$ [1].

\vspace{0.5em}\textbf{Q7(b).} $p = p_0 + \rho g h = 1.01\times10^5 + 1025 \times 9.81 \times 180$ [1]; $= 1.01\times10^5 + 1.81\times10^6 = \mathbf{1.91\times10^6\ \text{Pa}}$ [1].

\vspace{0.5em}\textbf{Q7(c).} $F = \rho g V = 1025 \times 9.81 \times 4.5\times10^{-3} = \mathbf{45.2\ \text{N}}$ [1]; weight of sphere $= 35 \times 9.81 = 343\ \text{N}$ [1]; $W > F$, so the sphere \textbf{sinks} [1].

\end{markscheme}

% ===== SECTION 9: REVISION CHECKLIST =====
\newpage
\section{Revision Checklist}

Use this checklist to track your understanding. Tick each box when you are confident:

\begin{tcolorbox}[colback=white, colframe=darkblue, breakable]
\renewcommand{\arraystretch}{1.5}
\begin{tabular}{@{} c p{10.5cm} r @{}}
\toprule
 & \textbf{Learning Objective} & \textbf{Confidence (1--3)} \\
\midrule
$\square$ & Explain what is meant by the centre of gravity of an object & \\
$\square$ & Define and calculate the moment of a force about a point & \\
$\square$ & Explain what a couple is and calculate the torque of a couple & \\
$\square$ & State and apply the principle of moments & \\
$\square$ & State the two conditions for equilibrium of a rigid body & \\
$\square$ & Use a vector triangle to represent three coplanar forces in equilibrium & \\
$\square$ & Define density and use \nf{\rho = m/V} & \\
$\square$ & Define pressure and use \nf{p = F/A} & \\
$\square$ & Derive and use the hydrostatic pressure equation \nf{\Delta p = \rho g \Delta h} & \\
$\square$ & Explain the origin of upthrust in terms of pressure difference & \\
$\square$ & Calculate upthrust using Archimedes' Principle \nf{F = \rho g V} & \\
$\square$ & Determine whether an object floats or sinks by comparing weight and upthrust & \\
\bottomrule
\end{tabular}
\vspace{0.5em}

\textit{Key: 1 = Need more work \quad 2 = Getting there \quad 3 = Confident}
\end{tcolorbox}

\vspace{1em}
\begin{motivationbox}
\centering
\textbf{\color{darkblue}Good luck with your revision!}\\[0.2em]
{\color{darkgray}\small Remember: always check \emph{both} equilibrium conditions, and use perpendicular distances when calculating moments. Drawing a clear force diagram before each problem will save you marks.}
\end{motivationbox}
